
% \vspace{-0.05in}
\insertToyResultsFigure
\vspace{-0.05in}
\vspace{-0.05in}

\section{Experiments}
\label{sec:experiments}

We compare \bocas to the following four baselines: 
(i) \gpucb,
(ii) the \gpeis criterion in \bayo~\citep{jones98expensive}, 
(iii) \mfgpucbones~\citep{kandasamy16mfbo} and
(iv) \mfsko, the multi-fidelity sequential kriging optimisation method
from~\citet{huang06mfKriging}.
All methods are based on GPs and we use the SE kernel for both the fidelity space
and domain.
The first two are not multi-fidelity methods, while the last two are finite multi-fidelity
methods\footnote{%
To our knowledge, the only other work that applies to continuous approximations
is~\citet{klein2015towards} which was developed specifically for hyper-parameter tuning.
Further, their implementation is not made available and is not straightforward to
implement.  }.
We have described the implementation details for all methods in
Appendix~\ref{sec:appImplementation}.


\subsection{Synthetic Experiments}
\label{sec:synthetic}


The results for the first set of synthetic experiments are given in Fig.~\ref{fig:toy}.
The title of each figure states the function used, and the dimensionalities $p,d$ of the
fidelity space and domain.
In all cases, the fidelity space was taken to be $\Zcal=[0,1]^p$ with
$\zhf = \one_p = [1,\dots,1]\in\RR^p$
being the most expensive fidelity.
For \mfgpucbones and \mfsko, we used $3$ fidelities ($2$ approximations) where the
approximations were obtained at $z=0.333\one_p$ and $z=0.667\one_p$ points in the fidelity
space.
To reflect the setting in our theory, we add Gaussian noise to the function value
when observing $g$ at any $(z,x)$.
This makes the problem more challenging than standard global optimisation problems
where function evaluations are not noisy.
The functions $g$, the cost functions $\cost$ and the noise variances $\eta^2$ are given
in Appendix~\ref{sec:appSynthetic}.

The first two figures in Fig.~\ref{fig:toy} are simple sanity checks.
In both cases, $\Zcal=[0,1]$ and $\Xcal=[0,1]$ and the functions were sampled from GPs. 
The GP was made known to all methods, i.e. the methods used the true GP in picking the
next point.
In the first figure, we used an SE kernel with bandwidth $0.1$ for $\phix$ and
$1.0$ for $\phiz$.
The large fidelity bandwidth causes $\gunc$ to be smooth across $\Zcal$ and
\mfgpucbtwos outperforms other baselines in this setting.
The curve starts mid-way in the figure as \bocas is yet to query at $\zhf$ up until
that point.
The second figure uses the same set up as the first except we used an SE
kernel with bandwidth $0.01$ for $\phiz$.
Even though $\gunc$ is highly  unsmooth across $\Zcal$,
\mfgpucbtwos does not perform poorly.
% It immediately skips the lower fidelities (curve starts early) as the threshold
% $\gamma(z)$ is large in~\eqref{eqn:candfidelst} and achieves as good regret
% as \gpucbs and \gpei.
This corroborates a claim that we made earlier that \bocas can naturally adapt
to the smoothness of the approximations.
The other multi-fidelity methods seem to suffer in this setting.



In the remaining experiments, we use some standard benchmarks for global optimisation.
We modify them to obtain $g$ and add noise to the observations.
As the kernel and other GP hyper-parameters are unknown, we learn them by maximising
the marginal likelihood every $25$ iterations.
This is a common heuristic used in the \bayo{} literature.
We outperform all methods on all problems
except in the case of the Borehole function where \mfgpucbones does better.
The last synthetic experiment is the Branin function given in Fig.~\ref{fig:branin}.
We used the same set up as above, but use 10 fidelities for \mfgpucbones and \mfskos
where the $k$\ssth fidelity is obtained at $z = \frac{k}{10}\one_p$ in the fidelity space.
Notice that the performance of finite fidelity methods deteriorate.
In particular, as \mfgpucbones does not share information across fidelities, the
approximations need to be designed carefully for the algorithm to work well.
Our more natural modelling assumptions prevent such pitfalls.
We next present two real examples in astrophysics and hyper-parameter tuning.
We do \emph{not} add noise to the observations, but treat it as optimisation tasks,
where the goal is to maximise the function.

\subsection{Astrophysical Maximum Likelihood Inference}

\label{sec:sn}


We use data on TypeIa supernova for maximum likelihood
inference on $3$ cosmological parameters, the Hubble constant $H_0\in(60,80)$, the
dark matter fraction $\Omega_M\in(0,1)$ and dark energy fraction $\Omega_\Lambda\in(0,1)$,
hence $d=3$.
The likelihood is given by the Robertson-Walker metric, the computation of which
requires a one dimensional numerical integration for each point in the dataset.
Unlike typical maximum likelihood problems, here the likelihood is only accessible
via point evaluations.
\insertRealResultsFigure
We use the dataset from~\citet{davis07supernovae} which has data on $192$ supernovae.
We construct a $p=2$ dimensional multi-fidelity problem where we can choose 
between data set size $N\in[50,192]$ and perform the integration on grids of size 
$G\in[10^2, 10^6]$ via the trapezoidal rule.
As the cost function for fidelity selection, we used $\cost(N,G) = NG$ as
the computation time is linear in both parameters.
Our goal is to maximise the average log likelihood at $\zhf=[192, 10^6]$.
For the finite fidelity methods we use three fidelities with the approximations available
at $z=[97,2.15\times 10^3]$ and $z=[145, 4.64\times 10^4]$ (which correspond to
$0.333\one_p$ and $0.667\one_p$ after rescaling as in Section~\ref{sec:synthetic}).
The results are given in Fig.~\ref{fig:sn} where we plot the maximum average log
likelihood against wall clock time as that is the cost in this experiment.
The plot includes the time taken by each method to tune the GPs and determine the next
points/fidelities for evaluation.
% \bocas outperforms other methods on this problem.


\subsection{Support Vector Classification with $20$ news groups}
\label{sec:news}

We use the $20$ news groups dataset~\cite{joachims1996probabilistic} in a text
classification task.
We obtain the bag of words representation for each document, convert them to tf-idf
features and feed them to a support vector classifier.
The goal is to tune the regularisation penalty and the temperature of the rbf kernel
both in the range $[10^{-2}, 10^3]$. Hence $d=2$.
The support vector implementation was taken from scikit-learn.
We set this up as a $2$ dimensional multi-fidelity problem where we can choose
a dataset size $N\in[5000, 15000]$ and the number of training iterations
$T\in[20, 100]$.
Each evaluation takes the given dataset of size $N$ and splits it up into $5$ to perform
$5$-fold cross validation.
As the cost function for fidelity selection, we used $\cost(N,T) = NT$ as
the training/validation complexity is linear in both parameters.
Our goal is to maximise the cross validation accuracy at $\zhf = [15000, 100]$.
For the finite fidelity methods we use three fidelities with the approximations available
at $z=[8333,47]$ and $z=[11667, 73]$.
The results are given in Fig.~\ref{fig:news} where we plot the average cross validation
accuracy against wall clock time.


% Our goal is to use tune the 

